\lysh{13055}{4}

\begin{alist}
\item Για κάθε $x \in \mathbb{R}$ έχουμε:
\[f(2 + x) = (2 + x)^2 - 4(2 + x) + 5 = x^2 + 4x + 4 - 4x - 8 + 5 = x^2 + 1\]
και
\[f(2 - x) = (2 - x)^2 - 4(2 - x) + 5 = x^2 - 4x + 4 + 4x - 8 + 5 = x^2 + 1\]
οπότε $f(2 + x) = f(2 - x)$.

\item Με $x = 1,52$ η ισότητα του ερωτήματος ($\alpha$) δίνει
$f(2 + 1,52) = f(2 - 1,52)$, οπότε $f(3,52) - f(0,48) = 0$ (1)
ενώ με $x = 1,48$ δίνει
$f(2 + 1,48) = f(2 - 1,48)$, οπότε $f(3,48) - f(0,52) = 0$ (2)
Με τη βοήθεια των (1) και (2) παίρνουμε:
\[A = f(3,52) - f(0,52) + f(3,48) - f(0,48) = f(3,52) - f(0,48) + f(3,48) - f(0,52) = 0 + 0 = 0.\]

\item Το πλήθος των κοινών σημείων της $C_f$ με την ευθεία είναι ίσο με το πλήθος των λύσεων
της εξίσωσης $f(x) = 2x + \beta$. Με $\beta = -5$ έχουμε:
\[f(x) = 2x - 5 \Leftrightarrow x^2 - 4x + 5 = 2x - 5 \Leftrightarrow x^2 - 6x + 10 = 0\]
που είναι αδύνατη αφού έχει διακρίνουσα $\Delta = 36 - 40 = -4 < 0$.
Άρα η $C_f$ δεν έχει κοινά σημεία με την ευθεία, όταν $\beta = -5$.

\item Η $C_f$ έχει ένα τουλάχιστον κοινό σημείο με την ευθεία $y = 2x + \beta$, μόνο όταν η εξίσωση
$f(x) = 2x + \beta$ έχει μια τουλάχιστον λύση. Είναι:
\[f(x) = 2x + \beta \Leftrightarrow x^2 - 4x + 5 = 2x + \beta \Leftrightarrow x^2 - 6x + 5 - \beta = 0\]
Η εξίσωση έχει μια τουλάχιστον λύση μόνο όταν η αντίστοιχη διακρίνουσα είναι μη αρνητική. Είναι:
\[\Delta \ge 0 \Leftrightarrow 36 - 4(5 - \beta) \ge 0 \Leftrightarrow 9 - 5 + \beta \ge 0 \Leftrightarrow \beta \ge -4\]
οπότε η ζητούμενη μικρότερη τιμή του $\beta$ είναι η τιμή $\beta = -4$.
\end{alist}

